\section{Discussion}
\label{sec:discussion}

\paragraph{Version-sensitivity}
An earlier campaign found Prisma (then 5.22, on an in-process Rust query engine) tying the native
driver on the deep fetch while drawing about five application cores. The \emph{identical} harness on
the re-frozen toolchain did not reproduce it: application CPU now stays near one core and Prisma sits
at mid-to-low throughput (Section~\ref{sec:results}). The non-reproduction is a fact, its cause a
hypothesis, since the re-freeze moved the whole toolchain at once, so the headline is
version-sensitive~\cite{vanzyl2009optimisations}.

\paragraph{Implications for practice}
RQ1 and RQ3 agree the selected-configuration difference is specific to the access \emph{pattern}:
the native-relative spread is widest on the deep fetch, where a layer must hydrate rows and resolve
associations rather than merely issue SQL (the object-relational impedance
mismatch~\cite{ireland2009impedance} made concrete), and narrowest on aggregation. Round-trip count
alone does not explain throughput: single-query MikroORM is slowest on the PostgreSQL deep fetch
while two-query Drizzle outruns four-query Prisma (Supplement Table~S2;
Table~\ref{tab:deep_fetch}). Result-graph materialization and the other bundled implementation
differences are plausible contributors, but the study does not isolate them.

\paragraph{Single-row inserts and backend-stack transfer}
Single-row insert behavior is both stack- and treatment-specific, and the durability and wait-event
controls reported there are consistent with a shared commit-path contribution, not with attributing
a performance floor or a fixed share to the DBMS (Section~\ref{sec:results}).

The second same-host campaign preserves PostgreSQL insert ranks but few MySQL positions
(Section~\ref{sec:results}), so we claim no durable fine ordering for MySQL insert. The recurring, material insert reversal
is narrower: Prisma exceeds MikroORM on PostgreSQL and trails it on MySQL in both
campaigns, with paired intervals excluding equality in opposite directions
(Supplement Table~S54).

Every promoted reversal involves Prisma, whose MySQL path uses the MariaDB adapter rather than
\texttt{mysql2}, so backend-stack sensitivity here is inseparable from that driver substitution and
describes one treatment, not ORMs as a class. \emph{Observed:} every reversal clearing the bar
involves the one layer whose MySQL driver differs. \emph{Interpretation:} at the stated bar only the
asymmetric treatment reverses durably, while whole-ordering agreement stays imperfect either way,
reported descriptively rather than as a further reversal claim. \emph{Not identified:} the design
cannot attribute the reversals to the DBMS rather than the adapter, the driver, or Prisma's
implementation, because holding out the treatment holds out the substitution; the hold-out followed
observing that all four involved one treatment, and its criterion is a design fact fixed
independently of the outcomes. This supports backend-stack sensitivity in the tested
implementations, leaving cross-host transfer unresolved. A five-replicate, one-layer-per-tier
six-statement transaction remains exploratory and does not extend the single-row-insert finding to
writes (Supplement Table~S8).

\paragraph{Methodological lessons: a checklist for access-layer benchmarking}
Several failures changed or could have changed the comparison; we distil them into a reusable
checklist (\texttt{notes/supplement-methods}): an \textbf{aggregation fan-out} and a driver
result-shape mismatch, both caught by specification and differential admission; \textbf{OFFSET
pagination} collapsing on MySQL, easily misread as an access-layer weakness; \textbf{write-state
contamination}, where MySQL allocated inserts below the cleanup floor; and a \textbf{query-formulation
confound} where a mis-scoped pre-aggregation re-scanned the whole table, a query-plan artifact
masquerading as an access-layer effect. The write path needs its own gate: MikroORM~7's dropped
\texttt{persistAndFlush} made every insert throw a 500 faster than a real one, briefly \emph{leading}
MySQL until the non-2xx counter flagged it. Because a 2xx alone would miss a \emph{silently}
incorrect write, the gate also validates post-write state. None is schema-specific; each extends the
database-benchmarking pitfalls of Fruth et al.~\cite{fruth2022tail} and Raasveldt et
al.~\cite{raasveldt2018fair} to the access-layer sub-genre.

\paragraph{Evidence for the protocol contribution}
Each stage had a concrete consequence: an independent seed specification caught shared state and
fixture defects that mutual equality would miss, treatment provenance made the loading-policy
sensitivity auditable, capacity characterization changed the tail interpretation, and the common-SQL
path narrowed but did not decompose the selected-treatment spread (Supplement Table~S50). Applied to
the nine external sources, the codebook yields claim licenses rather than a binary quality score
(Supplement Table~S37).

This demonstrates that the protocol is \emph{executable}: it admits and rejects concrete cells,
changed the reported interpretation at every stage exercised here, and codes an external corpus to a
reproducible per-stage result.
It does not demonstrate \emph{reusability by another person}, and two limits sharpen that: no source
was coded \emph{satisfied} on any stage, so the codebook's upper range is untested, and agreement
against this rater would measure agreement with the protocol's own designer.
The codebook, audit record, agreement scorer, and result-blind forms are released so the question can
be settled against this study.
The claim is therefore a \emph{candidate} discipline and reusable artifact, demonstrated once: not
invention of the constituent controls, not a complete theory of benchmarking validity, and not a
validated methodology.

\paragraph{Practical guidance}
The synthesis follows from RQ1--RQ3 and describes the benchmarked implementations in the tested
configuration (Section~\ref{sec:threats}), not access-layer categories in general.
Because the fan-out sweep varies breadth (0--500 children) but not depth, schema, or query mix, the
deep-fetch peak locates a difference among five probes on one synthetic schema: an observation to
re-measure on the target workload, not a law.
Where a hot path resembles the tested deep fetch the access layer is consequential and worth
benchmarking, the thin paths (native driver, Knex) leading here for throughput and for p99 at high
load, though at matched utilization the MySQL native driver falls to seventh of eight
(Section~\ref{sec:results}).
Where the tested workload is read-simple, or for the tested MySQL single-row insert, the observed
layer spread is smaller.
These are performance findings only: the choice also turns on qualities this study does not measure
(Section~\ref{sec:introduction}), which may outweigh them, and the common-SQL results are a compound
sensitivity contrast on the selected-configuration cost, not a licence to use raw SQL.
